\SourceLeaf{273}%
\noindent Wählte man hingegen den zehnmillionsten Theil des Erdquadranten selbst, nach obigen Dimensionen, zur Einheit, so würde sein
\[
\log A=6{,}8049902365
\]
\par\medskip\Needspace{4\baselineskip}\pdfbookmark[1]{Art. 6}{geodaesie1-art-6}\centerline{\textbf{6}}\nobreak\smallskip
\indent Die Berechnung der Breite auf der Kugel aus der Breite auf dem Ellipsoid kann füglich nach der Formel 3 geführt werden, wenn sie nur für wenige Fälle gefordert wird; für ausgedehntere Anwendungen hingegen wird der Gebrauch einer Reihe vortheilhaft sein. zu deren Entwickelung hier die nöthigen Formeln gegeben werden sollen.

Ich bezeichne eine unbestimmte Breite auf dem Ellipsoid, oder einen unbestimmten Werth von \(90^0-w\), durch \(P+p\), und die entsprechende Breite auf der Kugel, oder den Werth von \(90^0-U\) durch \(Q+q\). Nach dem \textsc{Taylor}schen Lehrsatze wird
\[
q=\frac{\mathrm{d}U}{\mathrm{d}w}\mathord{.}p-\frac{1}{2}\mathord{.}\frac{\mathrm{dd}U}{\mathrm{d}w^2}\mathord{.}pp+\frac{1}{6}\mathord{.}\frac{\mathrm{d}^3U}{\mathrm{d}w^3}\mathord{.}p^3-\frac{1}{24}\mathord{.}\frac{\mathrm{d}^4U}{\mathrm{d}w^4}\mathord{.}p^4+\text{u.\,s.\,w.}
\]
wo für die Differentialquotienten diejenigen bestimmten Werthe zu substituiren sind, welche zu \(p=0\), oder zu \(w=90^0-P\), \(U=90^0-Q\) gehören. Die successive Entwickelung der unbestimmten Differentialquotienten ergibt
\[
\frac{\mathrm{d}U}{\mathrm{d}w}=\frac{\alpha(1-ee)\sin U}{(1-ee\cos w^2)\sin w}
\]
\[
\begin{aligned}
\frac{\mathrm{dd}U}{\mathrm{d}w^2}={}&\frac{\alpha(1-ee)\sin U}{(1-ee\cos w^2)^2\sin w^2}\\
&\quad\{\alpha(1-ee)\cos U-\cos w+ee(\cos w^3-2\cos w\sin w^2)\}
\end{aligned}
\]
\[
\begin{aligned}
\frac{\mathrm{d}^3U}{\mathrm{d}w^3}={}&\frac{\alpha(1-ee)\sin U}{(1-ee\cos w^2)^3\sin w^3}\{\alpha\alpha(1-ee)^2(\cos U^2-\sin U^2)\\
&-3\alpha(1-ee)\cos U(\cos w-ee(\cos w^3-2\cos w\sin w^2))\\
&+2\cos w^2+\sin w^2-ee(4\cos w^4-2\sin w^4)\\
&+e^4(2\cos w^6-\cos w^4\sin w^2+6\cos w^2\sin w^4)\}
\end{aligned}
\]
Die beiden folgenden, welche ich gleichfalls entwickelt habe, setze ich um den Raum zu schonen in ihrer unbestimmten Form nicht hieher.

Die Substitution von \(w=90^0-P\), \(U=90^0-Q\) ergibt dann, wenn zugleich
\par\noindent anstatt \(\alpha\sin Q\) der Werth \(\sin P\) (nach Gleichung 10), und
\par\noindent anstatt \(\alpha\cos Q\) der Werth \(\frac{\cos P}{\cos\zeta\cos\eta}=\frac{\cos\theta\cos P}{\cos\varphi}\) (nach Gleichungen 18, 16, 23)
\par\noindent substituirt, und zur Abkürzung \(\cos P=c\), \(\sin P=s\) geschrieben wird,
